% sections/reference_graph.tex --- Reference Graph and Schema
\section{Reference Graph and Schema}
\label{sec:graph}

This section fixes the on-disk schema of the released benchmark.
The single artefact that a downstream user needs in order to train
and evaluate any of the baselines reported in
Section~\ref{sec:evaluation} is
\texttt{graphs/je\_network.csv}; companion files are available for
joining additional features back from the row-level journal-entry
table.

\subsection{File layout}
\label{sec:graph:layout}

The released benchmark archive contains the following top-level
files:

\begin{itemize}
  \item \texttt{graphs/je\_network.csv} --- the flat edge list,
    one row per debit-credit edge produced by the
    Section~\ref{sec:network} dispatch.  This is the primary input
    to all baselines.
  \item \texttt{journal\_entries.csv} --- the row-level
    $42$-column journal-entry table, one row per line item, used
    to recover entry-level features when needed.
  \item \texttt{coa.csv} --- the chart of accounts with ISO~21378
    \texttt{Type}, \texttt{Class}, and \texttt{Sub-Class} fields
    populated alongside the internal \texttt{AccountType}
    classification.
  \item \texttt{period\_close/trial\_balances/} --- per-period
    trial balances for spot-checking distributional tests.
  \item \texttt{labels/} --- per-entry anomaly labels
    (\texttt{anomaly\_labels.csv}, \texttt{fraud\_labels.csv},
    \texttt{quality\_labels.csv}); already projected onto the edge
    list, but exposed separately for users who wish to re-label.
\end{itemize}

The full set of $\sim$$60$ output directories produced by
\system{} is available in the same archive for users who require
non-anomaly features (master data, document flow, sub-ledger,
intercompany eliminations, etc.); these are documented in
\citep{ivertowski2026datasynth} and are not strictly required for
the baselines reported here.

\subsection{Edge-list schema}
\label{sec:graph:schema}

Table~\ref{tab:graph:schema} fixes the $13$ columns of
\texttt{graphs/je\_network.csv}.  All identifiers are UTF-8 strings;
all amounts are serialised as decimal strings (no IEEE 754) to
preserve exact precision; dates use ISO 8601 (\texttt{YYYY-MM-DD}).

\begin{table}[H]
  \centering
  \caption{Schema of \texttt{graphs/je\_network.csv}.}
  \label{tab:graph:schema}
  \small
  \begin{tabular}{p{1.6cm}p{2.5cm}p{8.0cm}}
    \toprule
    \textbf{Column} & \textbf{Type} & \textbf{Description}\\
    \midrule
    \texttt{edge\_id}            & UUID v5  &
      Deterministic identifier; UUID v5 of
      $(\text{document\_id}, \text{debit.line\_number},
        \text{credit.line\_number})$.\\
    \texttt{document\_id}        & string   &
      Source document identifier
      (e.g.\ \texttt{PO-2024-000123}).  \texttt{Manual} when no
      source document.\\
    \texttt{posting\_date}       & ISO date &
      Posting date of the underlying journal entry.\\
    \texttt{from\_account}       & string   &
      Debit (source) general-ledger account number.\\
    \texttt{to\_account}         & string   &
      Credit (destination) general-ledger account number.\\
    \texttt{from\_line\_id}      & UUID v5  &
      Stable line-level identifier of the debit line.\\
    \texttt{to\_line\_id}        & UUID v5  &
      Stable line-level identifier of the credit line.\\
    \texttt{amount}              & decimal  &
      Edge weight in the entry's currency, in absolute value.\\
    \texttt{confidence}          & decimal  &
      Per-entry confidence per Table~\ref{tab:confidence};
      $1.00$ for Methods A, B-distinct, C-distinct, and D.\\
    \texttt{predecessor\_edge\_id} & UUID v5 &
      Edge in the immediately preceding journal entry of the same
      document chain matching on \texttt{from\_account}; empty for
      chain heads and for journal entries outside any chain.\\
    \texttt{business\_process}   & enum     &
      One of \texttt{P2P}, \texttt{O2C}, \texttt{R2R},
      \texttt{H2R}, \texttt{A2R}, \texttt{S2C}, \texttt{MFG},
      \texttt{BANK}, \texttt{AUDIT}, \texttt{Treasury},
      \texttt{Tax}, \texttt{Intercompany}, \texttt{PROJECT},
      \texttt{ESG}.\\
    \texttt{is\_fraud}           & bool     &
      Edge inherits the \texttt{is\_fraud} flag of the underlying
      journal entry.\\
    \texttt{is\_anomaly}         & bool     &
      Edge inherits the broader \texttt{is\_anomaly} flag,
      covering fraud, error, process, statistical, and relational
      categories (Section~\ref{sec:datasynth:anomaly}).\\
    \bottomrule
  \end{tabular}
\end{table}

\subsection{Deterministic edge identity}
\label{sec:graph:determinism}

The \texttt{edge\_id} field is the UUID~v5 of the triple
$(\text{document\_id},\, \text{debit.line\_number},\,
  \text{credit.line\_number})$
under a fixed namespace seeded by \system{}'s
\texttt{DeterministicUuidFactory}.  Two consequences follow.
First, two runs of the generator with the same seed and
configuration produce identical \texttt{edge\_id} values for every
edge, and the released benchmark is therefore reproducible at the
edge identifier level.  Second, an edge identifier remains stable
under irrelevant rearrangements (re-ordering of unrelated postings,
re-numbering of unrelated documents) but \emph{changes} when any
of the three components changes; this is the property required
for cross-version differential analysis.

\subsection{Joinability}
\label{sec:graph:joinability}

The schema is joinable column-for-column against the row-level
journal-entry table via $(\text{document\_id},
  \text{from\_line\_id})$ and $(\text{document\_id},
  \text{to\_line\_id})$.  A typical query pattern, written in SQL
shorthand, is:

\begin{verbatim}
SELECT  e.*, j_d.account_class, j_d.account_sub_class,
              j_c.is_post_close
FROM    je_network e
LEFT JOIN journal_entries j_d
       ON j_d.line_id = e.from_line_id
LEFT JOIN journal_entries j_c
       ON j_c.line_id = e.to_line_id
WHERE   e.posting_date BETWEEN '2024-10-01' AND '2024-12-31';
\end{verbatim}

This recovers the full $42$-column entry context for any edge
without requiring duplication of features in the edge list.

\subsection{Predecessor chain}
\label{sec:graph:predecessor}

The \texttt{predecessor\_edge\_id} field traces booking chains across
documents.  For a payment posting that clears an open vendor
invoice, the predecessor edge is the corresponding
$\text{AP\_CONTROL} \to \text{VENDOR\_AUX}$ edge of the invoice
posting; for a customer receipt clearing an open AR balance, the
predecessor edge is the corresponding
$\text{AR\_CONTROL} \to \text{CUSTOMER\_AUX}$ edge.  Heads of chains
(purchase orders, sales orders, manual postings) carry an empty
\texttt{predecessor\_edge\_id}.  This field allows the construction
of \emph{document-chain features} (length of incoming chain,
expected vs.\ observed time between adjacent postings, etc.) without
requiring a separate document-flow export to be parsed.

\subsection{Alternate exports}
\label{sec:graph:alternate}

The same edge list is produced in three additional formats
for users who prefer typed graph artefacts:

\begin{itemize}
  \item PyTorch Geometric (\texttt{.pt}): node and edge tensors
    plus edge-level labels, ready to load into a
    \texttt{torch\_geometric.data.Data} object.
  \item Neo4j: CSV node and relationship files plus a Cypher
    \texttt{LOAD CSV} script.
  \item DGL: a serialised \texttt{DGLGraph} (binary) with edge
    labels.
\end{itemize}

These exports are produced from the same edge-list source and
share \texttt{edge\_id} values with the CSV file; switching format
does not alter any baseline numbers reported in
Section~\ref{sec:evaluation}.
